\documentclass[11pt]{article}

\usepackage[a4paper,margin=28mm]{geometry}
\usepackage[T1]{fontenc}
\usepackage[utf8]{inputenc}
\usepackage{lmodern}
\usepackage{microtype}
\usepackage{amsmath,amssymb,amsthm,mathtools}
\usepackage{booktabs}
\usepackage{enumitem}
\usepackage{xcolor}
\usepackage{hyperref}
\usepackage[nameinlink,noabbrev]{cleveref}

\hypersetup{
  pdftitle={Settlement-Oriented Computing: Foundational Skeleton},
  pdfauthor={Anthonie Reijm},
  colorlinks=true,
  linkcolor=black,
  citecolor=black,
  urlcolor=black
}

% ---------------------------------------------------------------------------
% Theorem environments
% ---------------------------------------------------------------------------

\newtheorem{definition}{Definition}[section]
\newtheorem{axiom}[definition]{Axiom}
\newtheorem{assumption}[definition]{Assumption}
\newtheorem{theorem}[definition]{Theorem}
\newtheorem{proposition}[definition]{Proposition}
\newtheorem{corollary}[definition]{Corollary}
\newtheorem{lemma}[definition]{Lemma}
\newtheorem{obligation}[definition]{Proof obligation}
\newtheorem{conjecture}[definition]{Conjecture}

\theoremstyle{remark}
\newtheorem{remark}[definition]{Remark}
\newtheorem{example}[definition]{Example}
\newtheorem{status}[definition]{Status}
\newtheorem{principle}[definition]{Principle}

% ---------------------------------------------------------------------------
% Notation
% ---------------------------------------------------------------------------

\newcommand{\Conf}{\mathcal{C}}
\newcommand{\Wit}{\mathcal{W}}
\newcommand{\Exec}{\mathcal{E}}
\newcommand{\Rel}{\mathbf{Rel}}
\newcommand{\World}{\mathbb{W}}
\newcommand{\U}{\mathbb{U}}
\newcommand{\Pres}{\mathcal{P}}
\newcommand{\Beh}{\mathcal{B}}
\newcommand{\id}{\mathrm{id}}
\newcommand{\powerset}{\mathcal{P}}
\newcommand{\dist}{\mathcal{D}}
\newcommand{\cand}{\operatorname{Cand}}
\newcommand{\adm}{\operatorname{Adm}}
\newcommand{\settle}{\operatorname{Settle}}
\newcommand{\beh}{\operatorname{beh}}
\newcommand{\enc}[1]{\ulcorner #1 \urcorner}
\newcommand{\weakstep}{\Rightarrow}
\newcommand{\realizes}[1]{\xRightarrow{#1}}
\newcommand{\carrier}[1]{\lvert #1 \rvert}

\newcommand{\Established}{\textsc{Established}}
\newcommand{\Conditional}{\textsc{Conditional}}
\newcommand{\Target}{\textsc{Target}}
\newcommand{\Open}{\textsc{Open}}

\title{\textbf{Settlement-Oriented Computing}\\
Foundational Skeleton for Executable Worlds}
\author{Anthonie Reijm}
\date{\today}

\begin{document}
\maketitle

\begin{abstract}
Settlement-Oriented Computing (SOC) is proposed as a computational formalism
for executable worlds. Its ontology contains only two primitive kinds:
\emph{configurations}, which are admissible world descriptions, and
\emph{witnesses}, which are typed lawful correspondences between
configurations. A configuration is not intrinsically a state, pattern, policy,
history, or program. Its interpretation depends on the witness in which it
participates. Consequently, SOC introduces no separate ontology of schemas and
instances.

This document is deliberately a foundational skeleton rather than a complete
paper. It fixes the minimal ontology, separates abstract witnesses from their
concrete realizations without creating a third primitive, places settlement
within a coalgebraic operational semantics, and states the universal-world
representation theorem that would internalize evaluation and yield reflection.
Each central claim is marked as established from definitions, conditional on
standard assumptions, or an open proof obligation.
\end{abstract}

\section{Scope and claim ledger}

The purpose of this document is to make the foundation falsifiable before a
full exposition is written. It therefore distinguishes three levels of claim:

\begin{center}
\begin{tabular}{@{}p{0.31\textwidth}p{0.17\textwidth}p{0.43\textwidth}@{}}
\toprule
Claim & Status & Required basis \\
\midrule
Configurations and witnesses form a category &
\Established &
Identity, typed composition, associativity, and unit laws. \\

Concrete witness realization respects composition &
\Conditional &
A relational interpretation preserving identities and composition. \\

Settlement defines a coalgebra &
\Established &
Once a successor functor and policy-induced transition map are fixed. \\

Behavior is determined by a final coalgebra &
\Conditional &
Existence of the relevant final coalgebra. \\

Policies and evaluation are internal to an SOC world &
\Target &
Construction of a universal world that simulates presented SOC worlds. \\

SOC is self-representationally closed &
\Target &
Derived from the universal-world representation theorem. \\

Programs and simulations require one ontology &
\Conditional &
Internal execution plus coalgebraic unfolding. \\

A concrete SOC realization is Turing complete &
\Open &
An explicit semantics-preserving encoding of a universal machine. \\
\bottomrule
\end{tabular}
\end{center}

SOC does not presently claim cartesian closedness, automatic termination,
confluence, efficient settlement, or a resolution of any complexity-theoretic
open problem.

\section{Minimal ontology}

\begin{definition}[Configuration]
A \emph{configuration} is any admissible description that may occur in an SOC
world. Let $\Conf$ denote the class of configurations.
\end{definition}

A configuration may be used as a concrete state, an intensional pattern, a
policy, a goal, an execution history, a program description, a proof state, or
an encoding of the settlement machinery. These are semantic roles, not
ontological subclasses.

\begin{definition}[Witness]
For configurations $A,B\in\Conf$, a \emph{witness}
\[
  w:A\longrightarrow B
\]
is a typed lawful correspondence from $A$ to $B$. Let
$\Wit(A,B)$ denote the class of witnesses from $A$ to $B$.
\end{definition}

The term ``witness'' includes static relations, observations, proofs,
reductions, state transitions, rewrites, and policy changes. Their operational
meaning is supplied by interpretation.

\begin{principle}[No schema ontology]
SOC does not distinguish schemas from instances as primitive kinds. A
configuration is interpreted intensionally or extensionally only relative to a
witness.
\end{principle}

\begin{remark}
This principle does not deny the semantic distinction between a reusable
description and one of its concrete applications. It denies that the
distinction requires a third primitive. Concrete application is represented by
a realization judgment, introduced next.
\end{remark}

\section{Witness-relative interpretation}

\begin{definition}[Relational interpretation]
Each witness $w:A\to B$ carries a realization relation
\[
  \rho_w \subseteq \Conf\times\Conf.
\]
We write
\[
  x \realizes{w} y
\]
when $(x,y)\in\rho_w$.
\end{definition}

The judgment $x\realizes{w}y$ means that the witness $w$, under its own
interpretation, realizes a concrete settlement possibility from $x$ to $y$.
It is a judgment about configurations and a witness, not a new ontological
entity.

\begin{definition}[Witness-relative source and target]
For $w:A\to B$, define
\[
  \operatorname{Src}(w)
  =
  \{x\in\Conf\mid \exists y\in\Conf,\;x\realizes{w}y\},
\]
and
\[
  \operatorname{Tgt}(w)
  =
  \{y\in\Conf\mid \exists x\in\Conf,\;x\realizes{w}y\}.
\]
\end{definition}

A single configuration may therefore be interpreted differently by distinct
witnesses. For example, one witness may require literal equality, another
structural compatibility, another logical entailment, and another a
probabilistic similarity threshold.

\begin{definition}[Object carrier]
The concrete carrier of a configuration $A$ is determined by its identity
witness:
\[
  \carrier{A}
  =
  \{x\in\Conf\mid x\realizes{\id_A}x\}.
\]
\end{definition}

\begin{assumption}[Typing of realizations]
For every witness $w:A\to B$,
\[
  \rho_w \subseteq \carrier{A}\times\carrier{B}.
\]
\end{assumption}

\begin{remark}[Rules and applications]
A witness may be reused across many concrete pairs $(x,y)\in\rho_w$. The
witness is the intensional lawful correspondence; the pair is one extensional
realization. Both $x$ and $y$ remain configurations, and no ``instance''
primitive is introduced.
\end{remark}

\section{Categorical kernel}

\begin{axiom}[Identity]
Every configuration $A$ has an identity witness
\[
  \id_A:A\to A.
\]
\end{axiom}

\begin{axiom}[Composition]
For witnesses $f:A\to B$ and $g:B\to C$, there is a composite witness
\[
  g\circ f:A\to C.
\]
\end{axiom}

\begin{axiom}[Associativity]
For every composable triple,
\[
  h\circ(g\circ f)=(h\circ g)\circ f.
\]
\end{axiom}

\begin{axiom}[Unit laws]
For every witness $f:A\to B$,
\[
  \id_B\circ f=f
  \qquad\text{and}\qquad
  f\circ\id_A=f.
\]
\end{axiom}

\begin{proposition}[Category formation]
Configurations and witnesses form a category, denoted
$\mathbf{SOC}$.
\end{proposition}

\begin{proof}
The objects are configurations, the morphisms are witnesses, and the preceding
axioms are exactly the category axioms.
\end{proof}

The proposition is foundational bookkeeping, not the main mathematical result.
The substantive obligation is compatibility between abstract witness
composition and concrete realization.

\begin{axiom}[Identity realization]
For every configuration $A$,
\[
  \rho_{\id_A}
  =
  \{(x,x)\mid x\in\carrier{A}\}.
\]
\end{axiom}

\begin{axiom}[Compositional realization]
For composable witnesses $f:A\to B$ and $g:B\to C$,
\[
  \rho_{g\circ f}
  =
  \rho_g\circ\rho_f,
\]
where the right-hand side is relational composition:
\[
  (x,z)\in\rho_g\circ\rho_f
  \iff
  \exists y,\;(x,y)\in\rho_f\land(y,z)\in\rho_g.
\]
\end{axiom}

\begin{proposition}[Realization functor]
Identity realization and compositional realization define a functorial
semantics from $\mathbf{SOC}$ into relations.
\end{proposition}

\begin{proof}
Identity witnesses are mapped to identity relations on their carriers, and
witness composition is mapped to relational composition. These are the
functor laws.
\end{proof}

\begin{corollary}[Sound path compression]
If
\[
  x_0\realizes{f_1}x_1
  \realizes{f_2}\cdots
  \realizes{f_n}x_n,
\]
then
\[
  x_0
  \realizes{f_n\circ\cdots\circ f_1}
  x_n.
\]
\end{corollary}

\begin{obligation}[Concrete BrixMS composition]
For a BrixMS realization, prove that the implementation's composition of
witnesses preserves the relational semantics above. In particular, matching,
binding, provenance, and update rules must not invalidate compositional
realization.
\end{obligation}

\section{Execution configurations and settlement}

\begin{assumption}[Internal packaging]
Finite tuples of configurations can be represented as configurations. In
particular, an execution configuration
\[
  e=\langle x,p,h\rangle\in\Exec\subseteq\Conf
\]
may contain a current world configuration $x$, a policy configuration $p$, and
a history configuration $h$.
\end{assumption}

This assumption is representational only. It does not yet prove that policy
evaluation is internally executable.

\begin{definition}[Candidate settlement]
For an execution configuration $e=\langle x,p,h\rangle$, a candidate
settlement is a tuple $(w,y,p',h')$ such that
\[
  x\realizes{w}y
\]
and the policy semantics admits the transition from $(p,h)$ to $(p',h')$.
The set of candidates is written $\cand(e)$.
\end{definition}

\begin{remark}[The evaluation boundary]
At this stage, the operation that interprets the policy configuration and
selects from $\cand(e)$ remains part of the metatheory. Encoding a policy as a
configuration does not by itself internalize evaluation. The universal-world
theorem in \cref{sec:universal} is the intended solution.
\end{remark}

\section{Coalgebraic settlement semantics}
\label{sec:coalgebra}

Let $\Exec$ be the class of execution configurations. A settlement discipline
chooses an endofunctor
\[
  F:\mathbf{Set}\to\mathbf{Set}
\]
describing the shape of possible next states.

\begin{definition}[Settlement coalgebra]
A settlement semantics is a coalgebra
\[
  \gamma:\Exec\longrightarrow F(\Exec).
\]
For each execution configuration, $\gamma$ returns the next-state structure
allowed by witness realization and policy.
\end{definition}

Typical choices include:
\[
\begin{array}{rcll}
F(X)&=&1+X, & \text{partial deterministic settlement},\\
F(X)&=&\powerset_{\!f}(X), & \text{finitely branching nondeterminism},\\
F(X)&=&\dist_{\!f}(X), & \text{finitely supported probability},\\
F(X)&=&K\times X, & \text{weighted or cost-annotated evolution}.
\end{array}
\]

More elaborate regimes may combine these constructions. The functor determines
the branching shape; the policy determines which witness realizations populate
that shape.

\begin{definition}[SOC program]
An SOC program is a pointed settlement coalgebra
\[
  (\Exec,\gamma,e_0),
\]
where $e_0$ is the initial execution configuration.
\end{definition}

\begin{definition}[Simulation]
A simulation is a finite or infinite path generated by repeated unfolding of
$\gamma$ from $e_0$, with branching resolved according to the chosen
settlement discipline.
\end{definition}

\begin{theorem}[Behavior map]
If $F$ has a final coalgebra
\[
  (\Beh,\zeta),
\]
then every SOC program admits a unique coalgebra morphism
\[
  \beh:\Exec\longrightarrow\Beh.
\]
\end{theorem}

\begin{proof}
This is the universal property of the final $F$-coalgebra.
\end{proof}

\begin{corollary}[Intensional and extensional views]
The pointed coalgebra $(\Exec,\gamma,e_0)$ is the intensional program
description, while $\beh(e_0)$ is its extensional behavior.
\end{corollary}

\begin{remark}
This is stronger and more precise than merely saying that a program and a
simulation are ``made of the same stuff.'' The coalgebra determines the
unfolding, and final semantics identifies behavior independently of one
particular presentation.
\end{remark}

\section{Finite presentations without schemas}

The universal construction cannot quantify over arbitrary proper-class-sized
worlds. It quantifies over effectively presented worlds.

\begin{definition}[Finite SOC presentation]
A finite SOC presentation is a finite syntactic datum
\[
  \Pres=(C_0,W_0,\mathsf{Real},\mathsf{Policy},e_0)
\]
where:
\begin{enumerate}[label=(\roman*)]
  \item $C_0$ is a finite set of configuration descriptions;
  \item $W_0$ is a finite set of witness descriptions;
  \item $\mathsf{Real}$ effectively enumerates the relation
        $\rho_w(x,-)$ for each described witness $w$ and configuration $x$;
  \item $\mathsf{Policy}$ effectively determines or enumerates admissible
        policy updates and selections;
  \item $e_0$ is an initial execution configuration.
\end{enumerate}
\end{definition}

The words ``description,'' ``pattern,'' and ``program'' indicate how a
configuration or witness is being interpreted. They do not introduce
additional primitive kinds. Open terms, ground terms, policies, and execution
states are all configurations.

\begin{assumption}[Effectiveness]
Configuration and witness descriptions are effectively enumerable; syntactic
equality is decidable; realization and admissibility are computably
enumerable.
\end{assumption}

\begin{definition}[Presented transition system]
The transition system generated by $\Pres$ has execution configurations as
states and a transition
\[
  e\longrightarrow_{\Pres}e'
\]
exactly when $\mathsf{Real}$ produces a witness realization compatible with
$e$ and $\mathsf{Policy}$ admits the resulting update.
\end{definition}

\section{The universal world}
\label{sec:universal}

\begin{definition}[Meta-configuration]
For a finite presentation $\Pres$ and execution configuration $e$, a
meta-configuration is an SOC configuration
\[
  \langle\enc{\Pres},\enc{e}\rangle.
\]
\end{definition}

\begin{definition}[Administrative and realizing steps]
A universal-world trace may contain:
\begin{enumerate}[label=(\roman*)]
  \item \emph{administrative steps}, which decode descriptions, enumerate
        realizations, and evaluate policy admissibility; and
  \item one \emph{realizing step}, which commits the represented settlement.
\end{enumerate}
Write
\[
  m\weakstep m'
\]
for zero or more administrative steps followed by one realizing step.
\end{definition}

\begin{conjecture}[Universal-world representation theorem]
There exists a finitely presented SOC world $\U$ such that, for every finite
SOC presentation $\Pres$ and execution configurations $e,e'$:
\begin{enumerate}[label=(\roman*)]
  \item \textbf{Soundness:}
  \[
    e\to_{\Pres}e'
    \implies
    \langle\enc{\Pres},\enc{e}\rangle
    \weakstep
    \langle\enc{\Pres},\enc{e'}\rangle.
  \]
  \item \textbf{Completeness:}
  every observable realizing step of $\U$ from
  $\langle\enc{\Pres},\enc{e}\rangle$ corresponds to a transition
  $e\to_{\Pres}e'$.
  \item \textbf{Progress:}
  if $e$ has an admissible represented successor, the universal world cannot
  diverge forever in administrative steps.
\end{enumerate}
Equivalently, the represented transition system and the corresponding
meta-level fragment of $\U$ are divergence-sensitive weakly bisimilar.
\end{conjecture}

\begin{status}
This is the principal target theorem. It replaces an internal-encoding axiom
with a construction obligation. The intended proof is a meta-interpreter
expressed as an SOC world.
\end{status}

\begin{obligation}[Construction of $\U$]
Specify a finite presentation for $\U$ and prove soundness, completeness, and
progress by induction over the meta-interpreter's administrative phases.
\end{obligation}

\begin{corollary}[Reflection, conditional]
If the universal-world representation theorem holds and $\U$ is itself
finitely presented, then $\U$ represents its own settlement semantics.
\end{corollary}

\begin{proof}
Instantiate the representation theorem with $\Pres=\Pres_{\U}$.
\end{proof}

\begin{corollary}[Self-representational closure, conditional]
Under the representation theorem, policies, policy revisions, scheduler state,
history, and evaluation can be represented and settled inside $\U$ without a
new ontology.
\end{corollary}

\begin{corollary}[Program--simulation co-location, conditional]
Under the representation theorem, a presented program and every represented
simulation of it occur inside the same universal SOC world: the program as a
meta-configuration with available observable steps, and the simulation as an
observable trace from that configuration.
\end{corollary}

\begin{remark}
The claim is co-location in one world, not set-theoretic identity between a
program description and one execution trace.
\end{remark}

\section{Behavioral adequacy}

Let $\gamma_{\Pres}$ be the coalgebra induced by a finite presentation and let
$\gamma_{\U}$ be the observable-step coalgebra of the universal world.

\begin{conjecture}[Adequacy]
Assuming the representation theorem and existence of the relevant final
coalgebra,
\[
  \beh_{\U}
  \bigl(\langle\enc{\Pres},\enc{e}\rangle\bigr)
  =
  \beh_{\Pres}(e).
\]
\end{conjecture}

\begin{proof}[Proof strategy]
Show that the representation relation is a divergence-sensitive weak
bisimulation over observable steps. Then apply the standard coalgebraic result
that bisimilar states have equal final behavior.
\end{proof}

\section{Worked kernel example: a two-counter machine}

A two-counter machine supplies a small, concrete test of the ontology and
settlement semantics.

\begin{definition}[Machine configuration]
A machine configuration is an SOC configuration
\[
  \langle q,m,n\rangle,
\]
where $q$ is an instruction label and $m,n\in\mathbb{N}$ are counter values.
\end{definition}

An increment instruction
\[
  q:\operatorname{INC}_1\to q'
\]
is represented by a witness $w_q$ with realization relation
\[
  \rho_{w_q}
  =
  \{
    (\langle q,m,n\rangle,\langle q',m+1,n\rangle)
    \mid m,n\in\mathbb{N}
  \}.
\]

A decrement-or-zero-test instruction
\[
  q:\operatorname{DECJZ}_1(q_{\mathrm{nz}},q_0)
\]
is represented by two witnesses:
\[
\begin{aligned}
  \rho_{w_q^{\mathrm{nz}}}
  &=
  \{
    (\langle q,m+1,n\rangle,
     \langle q_{\mathrm{nz}},m,n\rangle)
    \mid m,n\in\mathbb{N}
  \},\\
  \rho_{w_q^{0}}
  &=
  \{
    (\langle q,0,n\rangle,
     \langle q_0,0,n\rangle)
    \mid n\in\mathbb{N}
  \}.
\end{aligned}
\]

Instructions for the second counter are analogous.

\begin{proposition}[Single-step correspondence]
For the encoding above, a machine step
\[
  c\to_M c'
\]
exists if and only if there is an instruction witness $w$ such that
\[
  c\realizes{w}c'.
\]
\end{proposition}

\begin{proof}
By case analysis on the instruction at the control label of $c$. Each machine
instruction is represented by exactly the relation that defines its
operational step.
\end{proof}

\begin{obligation}[Universality transfer]
Complete the encoding with halting and control policy, then prove trace
correspondence. Combined with a standard universality theorem for two-counter
machines, this yields Turing completeness for any SOC realization capable of
executing the encoding.
\end{obligation}

\section{Position relative to adjacent formalisms}

Rewriting logic is the closest precedent for the universal-world target:
reflective rewriting admits a finitely presented universal theory, and Maude
implements meta-level computation. Universal coalgebra supplies the behavioral
semantics, finality, and bisimulation. Reflective Lisp systems expose the
evaluation boundary and reflective-tower problem. TLA treats actions as
state-transition relations and behaviors as state sequences. Graph
transformation gives rigorous matching and rewriting semantics, while the
Chemical Abstract Machine treats computation as rule-governed evolution of a
single state medium
\cite{clavel1996reflection,rutten2000universal,smith1984reflection,
lamport1994tla,konig2018graph,berry1992cham}.

The proposed SOC contribution is not the invention of categories, rewriting,
coalgebras, or reflection. Its candidate contribution is the combination of:

\begin{enumerate}[label=(\roman*)]
  \item one ontology of configurations and witnesses;
  \item witness-relative interpretation, avoiding a schema/instance ontology;
  \item policy-driven settlement as the operational organizing principle;
  \item coalgebraic behavior as the primary semantics; and
  \item a universal executable world as the reflection target.
\end{enumerate}

Whether this combination is novel and useful remains an empirical and
comparative research question.

\section{Foundation checklist}

Before expanding this skeleton into a full paper, the following questions
should have explicit answers.

\begin{enumerate}[label=\textbf{F\arabic*.}]
  \item Is every BrixMS runtime entity representable as a configuration or
        witness without hidden host-language objects?
  \item Does concrete witness composition satisfy compositional realization?
  \item Is matching entirely witness-relative, or are global interpretation
        rules still required?
  \item What exact functor $F$ describes each supported settlement discipline?
  \item Which final coalgebras are required, and under what existence
        conditions?
  \item Can the BrixMS meta-interpreter be finitely presented inside BrixMS?
  \item Can soundness, completeness, and progress of the meta-interpreter be
        proved?
  \item Does reflective ascent preserve observable behavior?
  \item Is the two-counter-machine encoding executable without hidden
        unbounded host-language primitives?
  \item What operational cost is introduced by witness discovery, policy
        evaluation, and reflective settlement?
\end{enumerate}

\section{Conclusion}

The foundation proposed here contains only configurations and witnesses.
There is no separate schema ontology: a configuration's semantic role is
determined by the witness that interprets it. Concrete applications are
realization judgments, not new entities. The categorical kernel organizes
composition, the relational interpretation connects abstract witnesses to
concrete settlement steps, and coalgebra supplies the dynamics and behavioral
semantics.

The unresolved center of the theory is evaluation. Encoding policies as
configurations is not sufficient; a universal world must faithfully execute
presented worlds, including itself. Constructing that world and proving weak
bisimulation is therefore the principal next theorem. If successful,
self-representational closure and program--simulation co-location become
derived results rather than assumptions.

\bibliographystyle{plain}
\bibliography{references}

\end{document}
